\documentclass[11pt,a4paper]{article}
\usepackage[top=22mm,bottom=24mm,left=24mm,right=24mm]{geometry}
\usepackage{fontspec}
\setmainfont[Path=fonts/,BoldFont=NotoSansKR-Bold.ttf,ItalicFont=NotoSansKR-Regular.ttf]{NotoSansKR-Regular.ttf}
\setsansfont[Path=fonts/,BoldFont=NotoSansKR-Bold.ttf]{NotoSansKR-Regular.ttf}
\setmonofont{Latin Modern Mono}
\XeTeXlinebreaklocale "ko"
\XeTeXlinebreakskip=0pt plus 0.1em
\usepackage{amsmath,amssymb,amsthm}
\usepackage{booktabs,longtable,array}
\usepackage{enumitem}
\usepackage{xcolor}
\definecolor{inkblue}{HTML}{173B60}
\usepackage[colorlinks=true,linkcolor=inkblue,citecolor=inkblue,urlcolor=inkblue,unicode]{hyperref}
\usepackage{fancyhdr}
\usepackage{setspace}
\setstretch{1.19}
\setlength{\parindent}{1em}
\setlength{\parskip}{3pt}
\setlength{\headheight}{15pt}
\pagestyle{fancy}
\fancyhf{}
\fancyhead[L]{\small 콜라츠 패리티 역상의 무리성}
\fancyhead[R]{\small Lee HaJin · 연구논문 초안}
\fancyfoot[C]{\thepage}
\newtheoremstyle{koreanplain}{7pt}{7pt}{\normalfont}{}{\bfseries}{.}{0.45em}{}
\theoremstyle{koreanplain}
\newtheorem{theorem}{정리}[section]
\newtheorem{lemma}[theorem]{보조정리}
\newtheorem{proposition}[theorem]{명제}
\newtheorem{corollary}[theorem]{따름정리}
\newtheorem{remark}[theorem]{주석}
\renewcommand{\proofname}{증명}
\renewcommand{\refname}{참고문헌}
\renewcommand{\abstractname}{초록}
\newcommand{\Q}{\mathbb Q}
\newcommand{\Z}{\mathbb Z}
\newcommand{\N}{\mathbb N}
\newcommand{\vTwo}{v_2}
\newcommand{\len}[1]{\lvert #1\rvert}
\newcommand{\htQ}{\operatorname{H}}
\newcommand{\ordz}{\operatorname{ord}_z}
\newcommand{\repo}{\texttt{hajin5305/collatz-research}}
\hypersetup{pdftitle={이진 균일 치환 고정점의 블록 부호화에 대한 콜라츠 패리티 역상의 무리성},pdfauthor={Lee HaJin},pdfsubject={AI를 활용한 연구논문 초안; 외부 심사 및 전체 우선권 검토 전}}
\title{\vspace{-1em}\textbf{이진 균일 치환 고정점의 블록 부호화에 대한\\콜라츠 패리티 역상의 무리성}\\[0.65em]
\large Irrationality of Collatz Parity Inverses for Block Codings\\of Binary Uniform Fixed Points}
\author{Lee HaJin}
\date{연구논문 초안 v0.1 · 2026-10-06}

\begin{document}
\maketitle

\begin{abstract}
이진 균일 치환의 고정점을 유한 이진 블록으로 부호화하여 얻은 무한 단어가 유리수 초기값의 콜라츠 패리티열로 실현될 수 있는지를 연구한다. $q\ge2$인 이진 $q$-균일 치환 $\sigma$의 일방향 고정점 $t$와 비어 있지 않은 이진 단어 $U,V$에 대해, $\kappa(0)=U$, $\kappa(1)=V$로 정의한 부호화를 생각한다. 콜라츠 패리티 역함수 $\Phi$에 대하여 $\Phi(\kappa(t))$가 유리수일 필요충분조건은 $t$가 최종 주기적이거나 $UV=VU$인 것이다. 치환의 원시성, 두 출력 블록의 같은 길이 또는 같은 1의 개수는 가정하지 않는다. 증명은 확대 블록의 역아핀 표현을 두 변수 급수 $F(z,y)$로 정리하고, 이동하는 두 번째 평가점의 로그 높이가 첫 번째 척도보다 낮은 차수로 성장함을 이용한다. $z$-차수 6의 고정 Padé 보조식으로 처음 13개 계수인 0차부터 12차까지를 소거한 뒤, 실제 평가값의 비영성을 따로 증명하여 $2^{13}>3^8$의 엄격한 높이 모순을 얻는다. Period-doubling과 비원시 Cantor 치환을 예로 제시하고, 유한 정확 검산의 범위를 기록한다. 정리는 특정 기호 구조의 산술적 실현을 제한하며, 일반 콜라츠 추측이나 모든 자동 패리티열의 배제를 결론하지 않는다. 본 연구와 초안 작성에는 생성형 AI가 활용되었다.
\end{abstract}

\noindent\textbf{주제어:} 콜라츠 사상; 2-adic 정수; 패리티열; 균일 치환; 블록 부호화; Padé 근사; 무리성.

\subsection*{English abstract}
We study rational realizability of infinite parity words obtained by block coding fixed points of binary uniform substitutions. Let $t=\sigma(t)$ be a one-sided fixed point of a binary $q$-uniform substitution, $q\ge2$, and let $\kappa(0)=U$ and $\kappa(1)=V$ be nonempty binary words. For the inverse Collatz parity map $\Phi$, we prove that $\Phi(\kappa(t))\in\Q$ if and only if $t$ is ultimately periodic or $UV=VU$. Neither primitivity nor equality of the output lengths or one-counts is required. The proof expresses the inverse map at every substitution scale through a fixed bivariate series $F(z,y)$. The second evaluation coordinate has logarithmic height negligible relative to the principal block length. A fixed Padé pair of degree at most six cancels the first thirteen coefficients. A separate nonvanishing argument for the actual evaluations then yields incompatible divisibility and height estimates, with the strict gap $2^{13}>3^8$. Period-doubling and nonprimitive Cantor controls illustrate the scope. Exact finite checks supplement the written proof. The theorem does not assert transcendence throughout the family or exclude arbitrary automatic Collatz itineraries. Generative AI was used in the research and preparation of this draft; the disclosure appears in Section~\ref{sec:ai}.

\section{서론}

콜라츠 사상의 반복을 각 단계의 짝홀성으로 기록하면, 정수 궤도에 관한 질문을 무한 이진 단어의 산술적 실현 가능성에 관한 질문으로 바꿀 수 있다. 이 관점의 기본 도구는 2-adic 정수와 패리티열을 연결하는 켤레사상이다~\cite{BL1996}. 모든 무한 이진 단어에는 유일한 2-adic 초기값이 대응하지만, 그 초기값이 유리수 또는 양의 보통정수라는 조건은 추가적인 제약이다.

본 논문은 이진 균일 치환의 고정점과 두 유한 블록의 부호화로 만들어지는 패리티열을 다룬다. 주요 결과는 이 클래스에서 유리 초기값을 갖는 경우의 필요충분조건이다. 부호화 전의 제어 단어가 최종 주기적이거나 두 출력 블록이 가환하면 패리티열은 최종 주기적이고, 유리 초기값을 갖는다. 두 조건이 모두 실패하면 대응하는 2-adic 초기값은 무리수이다.

주요 기술적 지점은 두 확대 블록의 길이와 1의 개수가 서로 달라도 같은 무한 초기값을 모든 확대 수준에서 비교한다는 데 있다. 이 차이는 두 번째 변수 $y_k$로 모인다. 이진 균일 치환의 계수 행렬이 주는 보조 척도 $|\delta|^k$는 주 척도 $q^k$보다 느리게 성장하므로, 이동 평가점의 높이를 제어할 수 있다. 이를 고정된 Padé 보조식의 실제 비영성과 결합한다.

\subsection{선행연구와 기여의 위치}

Bernstein과 Lagarias~\cite{BL1996}의 패리티 켤레사상은 본문의 기본 설정이다. $p$-adic 급수의 Padé 근사와 높이를 이용하는 방법에는 오래된 선행이 있으며, Bugeaud와 Yao~\cite{BY2017}는 Mahler형 $p$-adic 값과 Thue--Morse 관련 값에 대한 결과를 제시한다. 따라서 Padé 근사나 2-adic 높이 비교 자체를 새로운 원리로 주장하지 않는다.

Thue--Morse 제어의 블록 부호화에는 더 직접적인 선행도 있다. Allikvere~\cite{Allikvere2026}의 미심사 원고는 같은 길이의 서로 다른 출력 블록에 대한 초월성을 다룬다. 임의의 비가환 $U,V$에 대해 $UV$와 $VU$로 재그룹화하면 이 특수 경우와의 중복이 생긴다. Sharpe~\cite{Sharpe2026}의 미심사 원고는 $1\mapsto110$, $0\mapsto011$이라는 특정 균일 치환의 패리티열과 그 이동을 배제하며, 반복 평가와 높이 비교를 사용한다. 이 문헌들의 특수 경우와 본문의 전체 이진 균일 클래스는 가정을 나누어 비교해야 한다.

일반적인 $p$-adic 자리수 전개에 관한 초월성 정리~\cite{AB2007}를 적용할 때에도, 궤도의 패리티 비트와 $\Phi$ 자체의 2-adic 자리수를 구별해야 한다. 본문 급수에는 누적 1의 개수에 따라 달라지는 3의 거듭제곱 분모가 들어간다. 다변수 $p$-adic Mahler 값에 대한 정리들도 존재하지만~\cite{Ide2019}, 그 정리의 허용성 및 평가 비소멸 가정을 본문의 모든 치환에 대해 확인하는 일은 별도 문제이다.

본문은 정확한 전체 분류와 자족적인 증명을 제시한다. 검토한 문헌만으로 이 전체 결과의 최초성을 확정하지 않는다. 저장소의 최신 문헌 감사에서도 전체 BU 정리의 외부 우선권은 미확인으로 남아 있다~\cite{Repo2026}. 다음 절부터의 증명은 선행 초월성 정리나 밀도 정리를 추가 가정으로 사용하지 않는다.

\section{정의와 표기}

$\N_0=\{0,1,2,\ldots\}$라 하고, $\Z_2$와 $\Q_2$를 각각 2-adic 정수환과 2-adic 수체라 한다. $\vTwo(2)=1$로 정규화하고 $\vTwo(0)=+\infty$로 둔다. 보통 유리수는 $\Q_2$ 안에 자연스럽게 포함된 것으로 본다.

본문의 콜라츠 사상은 다음의 단축 사상 $T:\Z_2\to\Z_2$이다.
\begin{equation}\label{eq:T}
T(x)=\begin{cases}
x/2,&x\equiv0\pmod2,\\
(3x+1)/2,&x\equiv1\pmod2.
\end{cases}
\end{equation}
패리티열은 $\pi(x)_j\equiv T^j(x)\pmod2$로 정의한다. 무한 이진 단어 $w=w_0w_1\cdots$에 대해 $S_j=\sum_{i<j}w_i$라 놓으면 역함수는
\begin{equation}\label{eq:Phi}
\Phi(w)=-\sum_{j\ge0}\frac{w_j2^j}{3^{S_{j+1}}}\in\Z_2
\end{equation}
로 주어진다. 각 항의 2-adic valuation이 $j$ 이상이므로 급수는 수렴한다. 첫 비트에 따른 직접 계산으로 $\Phi(0v)=2\Phi(v)$와 $\Phi(1v)=(2\Phi(v)-1)/3$을 얻으며, 따라서 $\pi(\Phi(w))=w$이다. 이 표기는 고전적 패리티 켤레사상과 일치한다~\cite{BL1996}.

유한 단어 $u$의 길이를 $N(u)$, 1의 개수를 $A(u)$라 쓴다. 빈 단어를 제외한 모든 이진 유한 단어의 집합은 $\{0,1\}^{+}$이다. 단어의 곱은 이어 붙이기를 뜻한다. $t$가 \emph{최종 주기적}이라는 것은 어떤 유한 접두 $p$와 비어 있지 않은 유한 단어 $r$에 대해 $t=prrr\cdots$로 쓸 수 있다는 뜻이다.

치환 $\sigma:\{0,1\}^{*}\to\{0,1\}^{*}$가 $q$-균일이라는 것은 $N(\sigma(0))=N(\sigma(1))=q$라는 뜻이다. $t=\sigma(t)$는 일방향 무한 고정점이다. $\kappa(0)=U$, $\kappa(1)=V$는 $U,V\ne\varnothing$인 비소거 부호화이다. $UV=VU$일 때 두 블록이 가환한다고 한다.

기약 유리수 $a/b$ ($b>0$)의 높이는
\begin{equation}\label{eq:height}
\htQ(a/b)=\max\{|a|,b\}
\end{equation}
로 둔다. 모든 $O(\cdot)$와 $o(\cdot)$는 치환과 부호화를 고정한 뒤 $k\to\infty$에 대한 것이다. 아래의 $L$은 고정된 양의 실수이며, 창의 길이 또는 별도의 자유 매개변수가 아니다.

\section{주정리}

\begin{theorem}[이진 균일 고정점의 유리성 분류]\label{thm:main}
$q\ge2$인 이진 $q$-균일 치환 $\sigma$와 일방향 고정점 $t=\sigma(t)$를 고정하자. $U,V\in\{0,1\}^{+}$에 대해 $\kappa(0)=U$, $\kappa(1)=V$로 두면
\begin{equation}\label{eq:main}
\Phi(\kappa(t))\in\Q
\quad\Longleftrightarrow\quad
\bigl(t\text{가 최종 주기적}\bigr)\ \text{또는}\ UV=VU.
\end{equation}
치환의 원시성, $N(U)=N(V)$ 또는 $A(U)=A(V)$는 가정하지 않는다.
\end{theorem}

\begin{corollary}[정수 실현의 배제와 유한 변경]\label{cor:integer}
정리~\ref{thm:main}의 가정에서 $t$가 최종 주기적이지 않고 $UV\ne VU$이면, 어떤 $x\in\Q\cap\Z_2$도 전체 패리티열 $\kappa(t)$를 갖지 않는다. 특히 양의 보통정수 초기값은 없다. 이 결론은 패리티열의 유한 이동 및 임의의 유한 패리티 접두 부착에도 성립한다.
\end{corollary}

정리는 \emph{제어 알파벳이 이진이고 치환이 균일한} 경우를 다룬다. 임의의 유한 알파벳에서 출발하는 모든 자동수열이나 일반 비균일 치환으로 범위를 넓히지는 않는다. 또한 주정리의 비유리 결론은 전체 클래스의 초월성을 뜻하지 않는다.

\section{유한 블록의 역아핀 구조}

유한 이진 단어 $u=u_0\cdots u_{N-1}$에 대하여
\begin{equation}\label{eq:B}
B_0=0,\qquad B_{j+1}=3^{u_j}B_j+u_j2^j,\qquad B(u)=B_N
\end{equation}
로 정의하고
\begin{equation}\label{eq:clambda}
c(u)=-\frac{B(u)}{3^{A(u)}},\qquad
\lambda(u)=\frac{2^{N(u)}}{3^{A(u)}}
\end{equation}
로 둔다.

\begin{lemma}[역아핀 항등식과 분자 상계]\label{lem:affine}
유한 단어 $u,v$와 무한 꼬리 $w$에 대해
\begin{align}
\Phi(uw)&=c(u)+\lambda(u)\Phi(w),\label{eq:affine}\\
c(uv)&=c(u)+\lambda(u)c(v),\label{eq:concat}\\
0\le B(u)&\le3^{N(u)}-2^{N(u)}<3^{N(u)}.\label{eq:Bbound}
\end{align}
또한 같은 길이의 서로 다른 단어 $u,v$가 처음으로 위치 $r$에서 다르면
\begin{equation}\label{eq:isometry}
\vTwo(c(u)-c(v))=r.
\end{equation}
\end{lemma}

\begin{proof}
식~\eqref{eq:affine}은 식~\eqref{eq:Phi}를 접두와 꼬리로 나누어 얻는다. 빈 꼬리 뒤에 $0^\infty$를 붙이면 식~\eqref{eq:concat}도 따른다. 재귀식~\eqref{eq:B}은 비트 하나를 0에서 1로 바꿀 때 이후 분자를 감소시키지 않는다. 같은 길이에서 모든 비트가 1인 경우가 최대이며, 그때 분자는 $3^N-2^N$이다.

식~\eqref{eq:isometry}에서는 공통 접두를 제거한다. 공통 접두의 기울기는 2-adic valuation이 정확히 $r$이다. 제거한 두 꼬리의 첫 비트가 서로 다르므로 두 $c$는 하나가 짝수, 다른 하나가 홀수인 2-adic 정수이다. 따라서 꼬리 상수의 차이의 valuation은 0이고, 원래 차이의 valuation은 $r$이다.
\end{proof}

\begin{lemma}[두 단어 부호화의 단사성]\label{lem:code}
$U,V\in\{0,1\}^{+}$일 때 $\kappa(0)=U$, $\kappa(1)=V$가 유한 단어에서 단사일 필요충분조건은 $UV\ne VU$이다.
\end{lemma}

\begin{proof}
비단사 관계에서 공통 첫 부호어를 제거하면 서로 다른 $U,V$로 시작하는 같은 출력이 남는다. 그러면 짧은 단어는 긴 단어의 접두이다. $U=V$이면 가환한다. 그렇지 않고 $V=UW$라 하자. 추상 부호화 $0\mapsto0$, $1\mapsto01$은 끝에서부터 유일하게 해독되므로 단사이다. 따라서 원래 비단사 관계는 $\{U,W\}$의 비단사 관계를 유도한다. 총 길이에 대한 귀납으로 $UW=WU$를 얻고, 따라서 $U,V$도 가환한다. 반대로 $UV=VU$이면 추상 단어 $01$과 $10$의 출력이 같으므로 비단사이다.
\end{proof}

이제부터 정리~\ref{thm:main}의 어려운 방향을 위해 $t$는 최종 주기적이지 않고 $UV\ne VU$라고 가정한다. $\sigma(0)=\sigma(1)$이면 고정점이 그 공통 블록의 반복이므로 이 경우는 없다. 서로 다른 같은 길이의 두 치환상은 단사 부호화를 이루며, 보조정리~\ref{lem:code}와 합성하면
\[
W_{b,k}=\kappa(\sigma^k(b))\quad(b=0,1)
\]
에 대해 $W_{0,k}W_{1,k}\ne W_{1,k}W_{0,k}$이다. 이 두 단어는 같은 길이이므로
\begin{equation}\label{eq:D}
\begin{split}
D_k&=c_{1,k}(1-\lambda_{0,k})-c_{0,k}(1-\lambda_{1,k})\\
&=c(W_{1,k}W_{0,k})-c(W_{0,k}W_{1,k})\ne0
\end{split}
\end{equation}
이다. 여기서 $c_{b,k}=c(W_{b,k})$, $\lambda_{b,k}=\lambda(W_{b,k})$이며, $N_{b,k},A_{b,k},B_{b,k}$도 같은 방식으로 정의한다.

\section{균일 치환의 두 척도와 고정 급수}

$a_b=A(\sigma(b))$, $\delta=a_1-a_0$라 하자. $\delta=q$이면 치환은 $0\mapsto0^q$, $1\mapsto1^q$이고 모든 일방향 고정점이 상수이다. $\delta=-q$이면 첫 글자가 바뀌므로 일방향 고정점이 없다. 따라서 현재의 비주기 가정 아래에서는
\begin{equation}\label{eq:delta}
|\delta|<q.
\end{equation}

\begin{lemma}[블록 성장과 낮은 차수의 차이]\label{lem:counts}
$n_0=N(U)$, $n_1=N(V)$, $\rho=a_0/(q-\delta)$,
$L=(1-\rho)n_0+\rho n_1>0$라 두면 $k\ge1$에서
\begin{align}
N_{0,k}&=Lq^k-\rho(n_1-n_0)\delta^k,\label{eq:N0}\\
N_{1,k}&=Lq^k+(1-\rho)(n_1-n_0)\delta^k,\label{eq:N1}\\
A_{1,k}-A_{0,k}&=(A(V)-A(U))\delta^k.\label{eq:Adecay}
\end{align}
특히 $N_{b,k}=Lq^k+o(q^k)$이다.
\end{lemma}

\begin{proof}
$f_k=A(\sigma^k(0))$라 놓으면 $f_0=0$이고
$f_{k+1}=a_0q^k+\delta f_k$이다. 따라서
$f_k=\rho(q^k-\delta^k)$이며,
$A(\sigma^k(1))=f_k+\delta^k$이다. 각 0과 1을 $U,V$로 바꾸어 길이를 합하면 식~\eqref{eq:N0}와 \eqref{eq:N1}을 얻고, 같은 계산을 1의 개수에 적용하면 식~\eqref{eq:Adecay}를 얻는다. $0\le\rho\le1$이고 $n_0,n_1>0$이므로 $L>0$이다.
\end{proof}

평가점을
\begin{equation}\label{eq:zy}
z_k=\lambda_{0,k},\qquad
y_k=\frac{\lambda_{1,k}}{\lambda_{0,k}}
=\left(\frac{2^{n_1-n_0}}{3^{A(V)-A(U)}}\right)^{\delta^k}
\end{equation}
로 정의한다. $y_k$는 양의 유리수이며 2-adic 단원일 필요는 없다. 보조정리~\ref{lem:counts}는
\begin{equation}\label{eq:smallheight}
\htQ(z_k)\le3^{N_{0,k}},\qquad
\log\htQ(y_k)=o(q^k),\qquad |\vTwo(y_k)|=o(q^k)
\end{equation}
를 준다.

고정된 제어 단어 $t$의 접두 1의 개수를 $s_n=\sum_{j<n}t_j$라 쓰고
\begin{equation}\label{eq:FG}
F(z,y)=\sum_{n\ge0}t_nz^ny^{s_n},\qquad
G(z,y)=\sum_{n\ge0}z^ny^{s_n}
\end{equation}
로 둔다. 이는 $\Z[y][[z]]$의 급수이며, 계수 비교로
\begin{equation}\label{eq:FGid}
(1-z)G(z,y)=1+z(y-1)F(z,y)
\end{equation}
가 성립한다. $(z_k,y_k)$에서 $n$번째 항의 valuation은
\[
nN_{0,k}+s_n(N_{1,k}-N_{0,k})\ge n\min(N_{0,k},N_{1,k})
\]
이므로 두 급수는 $\Q_2$에서 수렴한다.

$t=\sigma^k(t)$이므로 $w=\kappa(t)$는 모든 $k$에서 같은 제어 단어 $t$에 따라 $W_{0,k},W_{1,k}$를 이어 붙인 단어이다. $\xi=\Phi(w)$라 하면 역아핀 항등식과 식~\eqref{eq:FGid}는
\begin{equation}\label{eq:seedF}
\xi=\frac{c_{0,k}+D_kF(z_k,y_k)}{1-z_k}
\end{equation}
를 준다. $1-z_k\ne0$은 $2^{N_{0,k}}\ne3^{A_{0,k}}$에서 따른다. 특히 식~\eqref{eq:seedF}의 $\xi$는 $k$마다 달라지는 초기값이 아니다.

\section{유리 특수화와 고정 Padé 보조식}

\begin{lemma}[모든 비영 유리 특수화의 비유리성]\label{lem:NS}
최종 주기적이지 않은 이진 단어 $t$와 $\eta\in\Q\setminus\{0\}$에 대하여 $F(z,\eta)\notin\Q(z)$이다.
\end{lemma}

\begin{proof}
$\eta=1$에서 유리 급수의 계수는 충분히 뒤에서 고정 선형 점화식을 만족한다. 계수가 0 또는 1이므로 유한 개의 연속 비트 상태가 다음 비트를 결정하며, 따라서 최종 주기성이 따른다.

$\eta\ne1$에서 $F(z,\eta)$가 유리라면 식~\eqref{eq:FGid}에 의해 $G(z,\eta)$도 유리이다. $g_n=\eta^{s_n}\ne0$라 놓으면 충분히 큰 $n$에 대해 어떤 고정 $r\ge1$과 $c_i\in\Q$가 있어
\[
g_{n+r}=\sum_{i=0}^{r-1}c_i g_{n+i}
\]
이다. $g_n$으로 나누면 오른쪽은 길이 $r-1$의 비트 창 $t_n,\ldots,t_{n+r-2}$로 결정된다. 이 창은 $g_{n+r-1}/g_n$도 결정하므로, 다음 비율 $g_{n+r}/g_{n+r-1}$을 결정한다. 그 비율은 서로 다른 두 값 $1,\eta$ 중 하나이므로 다음 비트도 유일하게 결정된다. 유한 상태의 결정적 진화는 최종 주기성을 강제한다. $r=1$에서는 빈 상태가 다음 비트를 결정한다. 두 경우 모두 가정과 모순이다.
\end{proof}

\begin{lemma}[차수 6, 소거 차수 13의 보조식]\label{lem:Pade}
$P,Q\in\Z[y,z]$와 고정 정수 $D\ge0$가 존재하여
\begin{equation}\label{eq:pade}
\begin{gathered}
\deg_zP,\deg_zQ\le6,\qquad \deg_yP,\deg_yQ\le D,\\
E(z,y)=P(z,y)+Q(z,y)F(z,y)\in z^{13}\Z[y][[z]],
\end{gathered}
\end{equation}
이고, 모든 $\eta\in\Q\setminus\{0\}$에서 $E(z,\eta)$는 영 급수가 아니다. $E(z,y)=\sum_{n\ge0}e_n(y)z^n$라 쓰면 $\deg e_n\le D+n$이다.
\end{lemma}

\begin{proof}
$\Q(y)$ 위에서 $z$-차수 6 이하인 $P,Q$의 계수는 총 14개이다. $P+QF$의 0차부터 12차까지를 소거하는 조건은 13개의 동차 선형식이므로 비자명한 해가 있다. $Q=0$이면 $P=0$이 강제되므로 $Q\ne0$이다.

분모를 제거한 뒤 $P,Q$의 모든 $z$-계수의 공통 $\Q[y]$ 인자를 나누고 정수 배수를 취한다. 이에 따라 $P,Q\in\Z[y,z]$이고 그 계수들의 공통 인자는 상수뿐이다. 소거 항등식은 유지된다. 어떤 유리 $\eta$에서 두 특수화 다항식이 모두 0이면 모든 계수에 $y-\eta$가 나누어져 모순이다. $Q(z,\eta)=0$이면 $P(z,\eta)\ne0$이고, 그렇지 않으면 보조정리~\ref{lem:NS}에 의해 $P(z,\eta)+Q(z,\eta)F(z,\eta)\ne0$이다. 마지막 차수 상계는 $s_n\le n$과 유한 차수의 $P,Q$에서 직접 따른다.
\end{proof}

여기서 $D$는 클 수 있지만 $k$와 무관하게 한 번 정해진다. 유한 계산에서 특정 Padé 쌍을 얻는 일은 보조정리~\ref{lem:Pade}의 보편적 존재 증명과 구별된다.

\section{주정리의 증명}

\subsection{유리 초기값 가정에서 얻는 높이 상계}

$\xi=p/r\in\Q$, $r>0$라고 가정하자. 식~\eqref{eq:D}와 \eqref{eq:seedF}에서 $h_k=F(z_k,y_k)$도 유리수이며
\begin{equation}\label{eq:hk}
h_k=
\frac{3^{A_{1,k}}\bigl[(3^{A_{0,k}}-2^{N_{0,k}})p+B_{0,k}r\bigr]}
{r\bigl[B_{0,k}(3^{A_{1,k}}-2^{N_{1,k}})-B_{1,k}(3^{A_{0,k}}-2^{N_{0,k}})\bigr]}.
\end{equation}
분모는 $D_k\ne0$에 의해 0이 아니다. 식~\eqref{eq:Bbound}와 $A_{b,k}\le N_{b,k}$를 사용하면 어떤 $k$와 무관한 $C_\xi>0$에 대해
\begin{equation}\label{eq:hheight}
\htQ(h_k)\le C_\xi 3^{N_{0,k}+N_{1,k}}
\end{equation}
이다. 예를 들어 $C_\xi=4(|p|+r)$로 충분하다. 기약화는 이 상계를 증가시키지 않는다.

\subsection{실제 평가값의 비영성}

보조정리~\ref{lem:Pade}에서 $E$를 고정한다. 형식급수 $E\ne0$만으로는 $E(z_k,y_k)\ne0$이 따르지 않으므로, 이를 따로 보인다.

$y_k$가 유한 집합에 속하면 무한히 반복되는 값 $\eta\ne0$의 부분수열을 택한다. $E(z,\eta)$는 비영 급수이고 $\vTwo(z_k)=N_{0,k}\to\infty$이다. 첫 비영 차수를 $\ell$, $h=\max(0,-\vTwo(\eta))$라 놓으면 $n\ge\ell+1$에서
\[
\vTwo(e_n(\eta)z_k^n)\ge n(N_{0,k}-h)-Dh.
\]
이 꼬리 전체의 하계는 첫 항의 $\ell N_{0,k}+\vTwo(e_\ell(\eta))$보다 결국 커진다. 따라서 이 부분수열에서 충분히 큰 $k$에 대해 $E(z_k,\eta)\ne0$이다. $\delta=0,1,-1$ 또는 식~\eqref{eq:zy}의 밑이 1인 경우가 여기에 포함된다.

나머지 경우에는 $|\delta|\ge2$이고 식~\eqref{eq:zy}의 양의 유리 밑이 1이 아니다. 그러면 $y_k$는 서로 다른 값들이다. $E$의 첫 비영 계수 다항식을 $e_m(y)$라 하자. $m\ge13$이고, 비영 다항식은 유한 개의 근만 가지므로 충분히 큰 $k$에서 $e_m(y_k)\ne0$이다. 유리 다항식 평가의 분자 상계와 식~\eqref{eq:smallheight}로
\begin{equation}\label{eq:leadingheight}
\vTwo(e_m(y_k))=o(q^k)
\end{equation}
를 얻는다. 구체적으로 $d_m=\deg e_m$라 놓으면 다음의 양쪽 상계를 사용한다.
\[
-d_m\max(0,-\vTwo(y_k))
\le \vTwo(e_m(y_k))
\le \log_2\bigl(\|e_m\|_1\htQ(y_k)^{d_m}\bigr).
\]

$\eta_k^-=\max(0,-\vTwo(y_k))$라 놓으면 $n\ge m+1$에서
\begin{equation}\label{eq:tail}
\vTwo(e_n(y_k)z_k^n)
\ge n\min(N_{0,k},N_{1,k})-D\eta_k^-.
\end{equation}
꼬리 전체의 valuation 하계는 $(m+1)Lq^k+o(q^k)$인 반면, 첫 비영항의 valuation은 $mLq^k+o(q^k)$이다. 따라서 첫 항은 상쇄되지 않는다. 두 경우를 합하면 $E(z_k,y_k)\ne0$인 무한 부분수열을 얻는다.

\subsection{소거 차수와 유리 높이의 충돌}

식~\eqref{eq:pade}와 같은 꼬리 추정은 모든 $k$에서
\begin{equation}\label{eq:lower}
\vTwo(E(z_k,y_k))
\ge13\min(N_{0,k},N_{1,k})-D\eta_k^-
=13Lq^k+o(q^k)
\end{equation}
를 준다.

반면 $z_k=a/b$, $y_k=c/d$, $h_k=u/v$를 각각 기약형으로 쓰고 양의 분모를 택하면, $E=P(z_k,y_k)+Q(z_k,y_k)h_k$의 공통분모로 $b^6d^Dv$를 사용할 수 있다. 분모를 제거한 정수 분자를 $\mathcal N_k$라 하면
\begin{align}
|\mathcal N_k|
&\le C_{P,Q}\htQ(z_k)^6\htQ(y_k)^D\htQ(h_k)\notag\\
&\le C'_{P,Q,\xi}3^{7N_{0,k}+N_{1,k}}\htQ(y_k)^D.\label{eq:numheight}
\end{align}
실제 평가값이 비영인 부분수열에서 $\mathcal N_k\ne0$이다. 비영 유리수 $\mathcal N/\mathcal M$의 정수 분모 $\mathcal M>0$에 대하여
$\vTwo(\mathcal N/\mathcal M)\le\log_2|\mathcal N|$이므로
\begin{equation}\label{eq:upper}
\begin{split}
\vTwo(E(z_k,y_k))
&\le (7N_{0,k}+N_{1,k})\log_2 3
 +D\log_2\htQ(y_k)+O(1)\\
&=8Lq^k\log_2 3+o(q^k).
\end{split}
\end{equation}
그런데
\begin{equation}\label{eq:gap}
13>8\log_2 3
\quad\Longleftrightarrow\quad
2^{13}=8192>6561=3^8.
\end{equation}
따라서 식~\eqref{eq:lower}와 \eqref{eq:upper}는 양립할 수 없다. $\xi\in\Q$라는 가정이 모순이며, 어려운 방향의 증명이 끝난다.

이 높이 추정에서 차수 5와 11차 소거를 쓰면 $2^{11}<3^7$이므로 같은 모순을 얻지 못한다. 차수 6은 본 증명에 충분한 선택이며, 모든 가능한 증명 방법에서 최소라는 주장은 아니다.

\subsection{역방향과 따름정리}

$t$가 최종 주기적이면 비소거 부호화 $\kappa(t)$도 최종 주기적이다. 비어 있지 않은 주기 블록 $r$에 대해
\[
\Phi(r^\infty)=\frac{c(r)}{1-\lambda(r)}\in\Q,
\]
이며 $1-\lambda(r)\ne0$이다. 유한 접두는 식~\eqref{eq:affine}에 의해 유리성을 보존한다.

$UV=VU$이면 두 단어는 공통 비어 있지 않은 단어의 거듭제곱이다. 이 사실은 긴 단어에서 짧은 접두를 제거하는 길이 귀납으로 얻는다. 따라서 모든 $\kappa(t)$는 그 공통 단어의 무한 반복이며 유리 초기값을 갖는다. 주정리의 필요충분조건이 완성된다.

유한 접두를 붙이거나 제거한 두 초기값은 식~\eqref{eq:affine}의 비영 유리 기울기를 갖는 아핀 변환으로 연결된다. 따라서 무리성은 유한 접두 부착과 이동에서 보존되고, 따름정리~\ref{cor:integer}도 성립한다. \hfill$\square$

\section{예와 적용 범위}

\subsection{Period-doubling 제어}

$\sigma(0)=01$, $\sigma(1)=00$의 0에서 시작하는 고정점은
$t_n=\vTwo(n+1)\bmod2$이며, $q=2$, $\delta=-1$이다. 이 단어가 충분히 뒤에서 주기 $p=2^ab$ ($b$ 홀수)를 갖는다고 가정하자. 충분히 큰 $k>a$와 $n=2^k-1$을 택하면 $t_n=k\bmod2$이지만 $\vTwo(n+p+1)=a$이다. $k$의 짝홀성을 바꾸면 모순이므로 최종 주기적이지 않다.

$U=110$, $V=111$은 비가환이므로 부호화된 패리티 역상은 무리수이다. 치환상의 1의 개수는 각각 1과 0이므로, 이 예는 치환상의 같은 1의 개수를 가정하지 않는 범위를 보여준다.

\subsection{비원시 Cantor 제어와 이동 평가점}

$\sigma(0)=000$, $\sigma(1)=101$의 1에서 시작하는 고정점에서는 $q=3$, $\delta=2$이다. 길이 $3^k$의 접두에 1이 정확히 $2^k$개 있다. 1은 무한히 나타나지만 밀도는 0이다. 최종 주기적 이진 단어는 1이 유한 개이거나 그 밀도가 양수이므로 이 고정점은 최종 주기적이지 않다.

$U=110$, $V=111$을 택하면 $y_k=3^{-2^k}$이다. 따라서 유한한 특수화값만 고려하는 논증으로는 이 예를 처리할 수 없으며, 실제로 이동하는 $y_k$의 높이 제어가 사용된다. 이 치환은 0에서 1을 만들지 않으므로 원시적이지 않다.

같은 Cantor 제어에 $U=10$, $V=111$을 택할 수도 있다. 이제 출력 길이는 2와 3이고, 1의 개수는 1과 3이며
\[
y_k=(2/9)^{2^k}.
\]
$10111\ne11110$이므로 비가환이고, 주정리는 다시 무리성을 준다. 이는 주정리의 직접 적용 예이며 추가적인 독립 정리를 주장하는 것이 아니다.

\subsection{유한 실현과 무한 실현의 구별}

길이 $H$의 이진 패리티 접두는 법 $2^H$의 한 정수 잔여류로 실현된다. 따라서 임의로 긴 유한 접두마다 양의 정수 실현자를 찾을 수 있다는 사실은 주정리와 모순되지 않는다. 본문이 배제하는 것은 \emph{하나의 고정 유리 초기값}이 지정된 무한 단어 전체를 실현하는 경우이다. 실현자가 접두 길이에 따라 달라지는 유한 자료로 이 조건을 대체할 수 없다.

\section{정확 검산과 재현성}\label{sec:reproduce}

수학적 결론의 근거는 앞 절들의 서면 증명이다. 계산은 유한 블록 항등식, 구체적인 Padé 인증서 및 구현의 일치 여부를 점검한다. 본 초안 준비 중 2026-10-06에 기준 저장소의 임시 복사본에서 원 생성기, 별도 검사기와 독립 Padé 검산을 실행했다. Python 3.12.14와 SymPy 1.14.0을 사용했으며, 생성 인증서와 저장된 결과의 자료형을 포함한 JSON 값 비교는 모두 일치했다.

\begin{center}
\begin{tabular}{lr}
\toprule
정확 유한 검산 항목 & 재현 범위\\
\midrule
유한 블록 역아핀 항등식 & 254개 단어\\
첫 불일치 위치의 valuation & 10,795개 쌍\\
아핀 commutator와 가환 예외 & 900개 쌍\\
치환 계수 항등식 & 1,200개 경우\\
확대 블록 급수 항등식 & 144개 경우\\
Padé 인증서 & 6개\\
기존 회귀 검사 & 14개\\
\bottomrule
\end{tabular}
\end{center}

6개 Padé 예에서 Cantor의 첫 비영항은 14차이고, 나머지 다섯 예는 13차이다. 모두 보조정리~\ref{lem:Pade}의 $E\in z^{13}\Z[y][[z]]$ 조건과 일치한다. 별도 구현으로 재구성한 세 예의 첫 비영항은 다음과 같다. 인증서의 상수배 또는 다른 기저 선택에 따라 계수의 표현과 부호는 달라질 수 있다.

\begin{center}
\begin{tabular}{lcl}
\toprule
제어 단어 & 첫 비영 차수 & 첫 계수\\
\midrule
Period-doubling & 13 & $y^6+y^5+y^4$\\
Cantor & 14 & $y^6$\\
Thue--Morse & 13 & $-2y^9-3y^8-2y^7$\\
\bottomrule
\end{tabular}
\end{center}

본 연구의 기준 자료는 \repo 의 commit
\begin{center}\small\texttt{f751b9102355c516434ed9d0882bf4eb8608d95b}\end{center}
이다. 주정리 원문은 \texttt{binary\_uniform\_pade/THEORY\_KO.md}의 §§0-9이며, 저장소에 처음 기록된 날짜는 2026-10-02이다. 2026-10-04의 별도 검토와 2026-10-06까지의 문헌 감사가 이어진다. 위 commit은 초안의 고정 기준점이며 최초 발견일을 뜻하지 않는다.

제공되는 부록 묶음에는 원 생성기와 검사기, 유한 인증서, 기대 결과 및 이번 재현 요약을 담았다. 저장소는 작성 시점에 비공개이므로 공개 데이터 접근성을 이미 확보한 것으로 주장하지 않는다. 정식 제출 시에는 인용 가능한 보존본과 공개 범위를 저자가 확정해야 한다. 증명보조기를 이용한 형식 검증이나 외부 동료심사는 이번 검산에 포함되지 않는다.

\section{논의와 남은 문제}

정리~\ref{thm:main}은 지정된 치환 및 부호화 클래스에서 유리 초기값을 정확히 분류한다. 제어 단어의 비주기성과 출력 블록의 비가환성이 모든 확대 수준에서 비영 commutator를 보장하고, 이진 균일 치환의 두 척도가 이동 평가점을 제어한다. 이 결합이 본문 증명의 중심이다.

첫 번째 남은 문제는 전체 정리의 문헌상 위치를 확정하는 것이다. 단일 변수 Mahler형 선행과 Thue--Morse 특수 경우는 명확히 인용할 수 있지만, 본문의 모든 비원시·불균형 치환과 서로 다른 출력 길이를 기존 다변수 $p$-adic 정리의 가정에 넣는 완전한 포함 증명은 여기서 제공하지 않는다. 반대로 그런 포함이 불가능하다는 주장도 하지 않는다.

두 번째는 적용 클래스의 확장이다. 제어 알파벳을 늘리면 하나의 $F(z,y)$와 유리 특수화 보조정리로 모든 자유도를 다루기 어려워진다. 비균일 치환에서는 두 확대 블록이 공통 주항 $Lq^k$를 갖는다는 현재의 길이 제어를 다시 검토해야 한다. 이러한 확장은 추가 가정과 별도의 증명을 요구한다.

세 번째는 정량화이다. 현재 증명은 무한 패리티열의 유리 실현을 배제하지만, 임의의 유리 초기값 높이에 따른 유한 일치 길이의 명시적 최적 상계를 제시하지 않는다. Padé 보조식의 계수, 특수화 비영성의 시작점과 성장 오차를 유효하게 추적하면 이 문제를 연구할 수 있다.

마지막으로, 임의의 양의 정수 궤도가 본문의 정확한 치환 구조를 가져야 한다는 정리는 없다. 주정리는 일반 발산 배제나 비자명한 양의 정수 주기 배제를 제공하지 않으며, 콜라츠 추측은 미해결로 남는다.

\section{AI 활용 및 저자 책임 고지}\label{sec:ai}

본 연구와 원고 준비에는 생성형 AI가 활용되었다. 기반 저장소는 AI 보조 연구 기록을 포함하며, 본 초안 작성에는 OpenAI의 ChatGPT를 사용하였다. AI는 문헌 후보 탐색과 정리, 증명 구조의 초안 작성 및 비판적 검토, 계산 코드 작성과 검산 보조, 원고 구성과 표현 정리에 사용되었다. 개별 과거 연구 세션의 모든 모델명과 버전은 이 초안에서 확인하지 못했으므로 임의로 기재하지 않는다.

이 원고의 저자는 \textbf{Lee HaJin}이다. AI는 저자 또는 독립된 동료심사자로 기재하지 않는다. AI에 의한 자연어 검토와 유한 계산의 성공은 외부 전문가 심사나 증명보조기 검증을 의미하지 않는다. 본 문서는 저자의 최종 검토 전 초안이며, 제출 원고의 정확성, 인용, 독창성 및 공개 범위에 대한 최종 확인과 책임은 저자에게 있다. 저자 소속, 연락처, 연구비 및 이해상충에 관한 확인되지 않은 정보는 기재하지 않았다.

\appendix
\section{증명의 의존성과 보조자료 안내}

주정리의 어려운 방향은 다음의 유한한 논리 의존관계를 가진다.
\begin{enumerate}[leftmargin=2em]
\item 식~\eqref{eq:Phi}와 유한 블록 재귀식에서 역아핀 항등식 및 첫 불일치 valuation을 얻는다.
\item 두 단어 부호화의 단사성으로 모든 확대 수준의 $D_k\ne0$을 얻는다.
\item 이진 균일 치환의 계수 재귀식으로 $N_{b,k}=Lq^k+o(q^k)$와 $\log\htQ(y_k)=o(q^k)$를 얻는다.
\item 같은 초기값 $\xi$를 고정 급수의 값 $F(z_k,y_k)$와 연결한다.
\item 유리 특수화의 비유리성과 선형대수로 고정 Padé 보조식을 선택한다.
\item 실제 평가값의 비영성을 확보한 뒤, 13차 소거와 공통분모 높이 상계를 비교한다.
\end{enumerate}
이 경로에 유한 탐색의 완전성, 미증명 밀도 가설 또는 기존 초월성 정리의 새로운 적용은 필요하지 않다.

원본 기준점과 부록 파일의 대응은 부록 묶음의 README에 기록했다. 재현 요약은 별도 JSON 파일에 담았다.
\begin{itemize}[leftmargin=2em]
\item 파일 대응과 실행 안내: \path{supplement/README.txt}
\item 이번 재현의 범위와 일치 결과: \path{supplement/validation_summary.json}
\end{itemize}
독립 보조 스크립트가 함께 계산하는 실수 pseudo-trajectory 진단은 주정리의 증명 입력이 아니므로, 본문의 결과 표에서는 제외하였다.

\begin{thebibliography}{99}

\bibitem{BL1996}
D. J. Bernstein and J. C. Lagarias,
\emph{The $3x+1$ Conjugacy Map},
Canadian Journal of Mathematics \textbf{48}(6) (1996), 1154-1169.
\href{https://doi.org/10.4153/CJM-1996-060-x}{doi:10.4153/CJM-1996-060-x}.

\bibitem{BY2017}
Y. Bugeaud and J.-Y. Yao,
\emph{Hankel determinants, Padé approximations, and irrationality exponents for $p$-adic numbers},
Annali di Matematica Pura ed Applicata \textbf{196} (2017), 929-946.
\href{https://doi.org/10.1007/s10231-016-0602-7}{doi:10.1007/s10231-016-0602-7}.
선택 대조 위치: 저자 원고 p.13의 유리 평가점 Remark.

\bibitem{AB2007}
B. Adamczewski and Y. Bugeaud,
\emph{On the complexity of algebraic numbers. I. Expansions in integer bases},
Annals of Mathematics \textbf{165}(2) (2007), 547-565.
\href{https://doi.org/10.4007/annals.2007.165.547}{doi:10.4007/annals.2007.165.547}.
선택 대조 위치: Theorem 6.

\bibitem{Allikvere2026}
J. Allikvere,
\emph{Periodic shadows, drift pressure, and transcendence of the $3x+1$ conjugacy map},
미심사 프리프린트 (2026), 저장소 문헌 감사가 대조한 판본.
\href{https://doi.org/10.5281/zenodo.21543841}{doi:10.5281/zenodo.21543841}.
선택 대조 위치: Theorem 5.9.

\bibitem{Sharpe2026}
M. Sharpe,
\emph{Rung One of the Word-Complexity Ladder: a Machine-Verified Exclusion of an Automatic Collatz Itinerary by Mahler's Method in Degree One},
미심사 원고, 2026-08-28.
\href{https://github.com/msharpe248/collatz/blob/a2aa8a0dda45290fd01c36c253b4b4e14a7e812b/paper/rung1.tex}{고정 원문: commit a2aa8a0, paper/rung1.tex}.
선택 대조 위치: Theorems 1.1-1.2와 적용 범위 논의.

\bibitem{Ide2019}
H. Ide,
\emph{Algebraic independence of certain entire functions of two variables generated by linear recurrences},
arXiv:1908.06289, v1 (2019).
\href{https://arxiv.org/abs/1908.06289}{arxiv.org/abs/1908.06289}.
선택 대조 위치: §3.1의 조건과 §3.3의 Theorem 9.

\bibitem{Repo2026}
Lee HaJin,
\emph{Collatz Research: binary-uniform Padé theorem, exact certificates and literature audits},
연구 저장소 \repo, commit \texttt{f751b9102355c516434ed9d0882bf4eb8608d95b}, 2026-10-06 확인.
\href{https://github.com/hajin5305/collatz-research/tree/f751b9102355c516434ed9d0882bf4eb8608d95b}{고정 저장소 기준점}.
원문: \texttt{docs/research\_records/2026-10-02/binary\_uniform\_pade/THEORY\_KO.md}.
독립 검토: 2026-10-04 \texttt{literature\_value\_audit/uniform\_review}.
문헌 감사: 2026-10-06 \texttt{novelty\_frontier} 및 \texttt{novelty\_access\_followup}.

\end{thebibliography}
\end{document}
